% ============================================================
%  preamble.tex — styling, environments, macros
%  Edit this file to change the look of the whole book.
% ============================================================

% ---------- Page geometry ----------
\usepackage[
  paperwidth=17cm, paperheight=24cm,
  inner=2.2cm, outer=1.8cm, top=2.2cm, bottom=2.4cm,
  headsep=0.7cm, footskip=1.2cm
]{geometry}

% ---------- Encoding & fonts ----------
\usepackage[T1]{fontenc}
\usepackage[utf8]{inputenc}
\IfFileExists{lmodern.sty}{\usepackage{lmodern}}{}
% Nicer fonts if available on your machine (MiKTeX / full TeX Live).
% Harmless no-ops if the packages are missing.
\IfFileExists{newtxtext.sty}{\usepackage{newtxtext}}{}
\IfFileExists{newtxmath.sty}{\usepackage{newtxmath}}{}
\IfFileExists{inconsolata.sty}{\usepackage[varqu,scaled=0.95]{inconsolata}}{}
% Font expansion needs scalable fonts; fall back gracefully if absent.
\IfFileExists{lmodern.sty}{\usepackage{microtype}}{\usepackage[expansion=false]{microtype}}

% ---------- Core packages ----------
\usepackage{graphicx}
\usepackage{xcolor}
\usepackage{booktabs}
\usepackage{tabularx}
\usepackage{longtable}
\usepackage{enumitem}
\usepackage{listings}
\usepackage[most]{tcolorbox}
\usepackage{fancyhdr}
\usepackage{titlesec}
\usepackage{caption}
\usepackage{etoolbox}
\usepackage{imakeidx}
\usepackage[hidelinks]{hyperref}

\hypersetup{
  colorlinks=true,
  linkcolor=mafblue,
  urlcolor=mafblue,
  citecolor=mafblue,
  pdftitle={Microsoft Agent Framework for .NET Engineers},
  pdfauthor={Konrad Cinkusz}
}

% ---------- Palette ----------
\definecolor{mafblue}{HTML}{1F4E79}
\definecolor{mafteal}{HTML}{0E7C7B}
\definecolor{mafamber}{HTML}{B26A00}
\definecolor{mafred}{HTML}{9B2C2C}
\definecolor{mafgrey}{HTML}{5A5A5A}
\definecolor{maflight}{HTML}{F4F6F8}
\definecolor{codebg}{HTML}{FAFAFA}
\definecolor{codeframe}{HTML}{D8DEE4}
\definecolor{kw}{HTML}{0B5394}
\definecolor{str}{HTML}{9B2C2C}
\definecolor{cmt}{HTML}{4F7A4F}
\definecolor{typ}{HTML}{0E7C7B}

% ---------- Headings ----------
\titleformat{\chapter}[display]
  {\normalfont\huge\bfseries\color{mafblue}}
  {\normalfont\normalsize\color{mafgrey}\MakeUppercase{\chaptertitlename}\ \thechapter}
  {8pt}{\Huge}
\titlespacing*{\chapter}{0pt}{10pt}{28pt}
\titleformat{\section}{\normalfont\Large\bfseries\color{mafblue}}{\thesection}{0.8em}{}
\titleformat{\subsection}{\normalfont\large\bfseries\color{mafgrey}}{\thesubsection}{0.7em}{}

\pagestyle{fancy}
\fancyhf{}
\fancyhead[LE]{\small\nouppercase{\leftmark}}
\fancyhead[RO]{\small\nouppercase{\rightmark}}
\fancyfoot[LE,RO]{\small\thepage}
\renewcommand{\headrulewidth}{0.4pt}

% ============================================================
%  CODE LISTINGS
% ============================================================

\lstdefinelanguage{csharpx}{
  language=[Sharp]C,
  morekeywords={
    var,await,async,record,init,with,nameof,when,
    global,using,namespace,file,required,scoped,
    AIAgent,IChatClient,ChatOptions,AgentThread,Workflow,
    WorkflowBuilder,AgentWorkflowBuilder,AIFunctionFactory,ValueTask
  },
  sensitive=true
}

\lstdefinestyle{mafcode}{
  basicstyle=\ttfamily\footnotesize,
  keywordstyle=\color{kw}\bfseries,
  stringstyle=\color{str},
  commentstyle=\color{cmt}\itshape,
  numbers=left,
  numberstyle=\ttfamily\tiny\color{mafgrey},
  numbersep=8pt,
  backgroundcolor=\color{codebg},
  frame=single,
  rulecolor=\color{codeframe},
  framesep=6pt,
  breaklines=true,
  breakatwhitespace=false,
  postbreak=\mbox{\textcolor{mafgrey}{$\hookrightarrow$}\space},
  showstringspaces=false,
  tabsize=4,
  columns=fullflexible,
  keepspaces=true,
  captionpos=b,
  aboveskip=1.2em,
  belowskip=1.2em,
  literate={ł}{{\l}}1 {Ł}{{\L}}1 {ą}{{\k a}}1 {ę}{{\k e}}1
           {—}{{-{}-}}2 {–}{{-}}1 {‘}{{'}}1 {’}{{'}}1
           {“}{{"}}1 {”}{{"}}1 {°}{{\textdegree}}1 { }{{~}}1
}

\lstset{style=mafcode}

% Inline C# environment.  Usage:
%   \begin{csharp}[caption={Hello agent},label={lst:hello}]
%   ...code...
%   \end{csharp}
\lstnewenvironment{csharp}[1][]
  {\lstset{language=csharpx,style=mafcode,#1}}
  {}

\lstnewenvironment{shellcmd}[1][]
  {\lstset{language=bash,style=mafcode,numbers=none,keywordstyle=\color{mafteal}\bfseries,#1}}
  {}

\lstnewenvironment{xmlcode}[1][]
  {\lstset{language=XML,style=mafcode,#1}}
  {}

\lstnewenvironment{yamlcode}[1][]
  {\lstset{style=mafcode,basicstyle=\ttfamily\footnotesize,#1}}
  {}

% External source file — the preferred form once code lives in code/.
% Usage: \csfile{code/ch02/HelloAgent/Program.cs}{Caption}{lst:label}
\newcommand{\csfile}[3]{%
  \lstinputlisting[language=csharpx,style=mafcode,caption={#2},label={#3}]{#1}%
}

% Inline literals
\newcommand{\code}[1]{\texttt{\small #1}}
\newcommand{\pkg}[1]{\texttt{\small\color{mafblue}#1}}
\newcommand{\api}[1]{\texttt{\small\color{mafteal}#1}}

% ============================================================
%  ADMONITIONS
% ============================================================

\newtcolorbox{note}[1][]{
  enhanced, breakable,
  colback=maflight, colframe=mafblue,
  boxrule=0pt, leftrule=3pt,
  arc=0pt, outer arc=0pt,
  left=8pt, right=8pt, top=6pt, bottom=6pt,
  fonttitle=\bfseries\small,
  title={Note}, #1
}

\newtcolorbox{warning}[1][]{
  enhanced, breakable,
  colback=maflight, colframe=mafred,
  boxrule=0pt, leftrule=3pt,
  arc=0pt, outer arc=0pt,
  left=8pt, right=8pt, top=6pt, bottom=6pt,
  fonttitle=\bfseries\small,
  title={Watch out}, #1
}

% For anything that is preview / alpha / likely to move.
\newtcolorbox{versionbox}[1][]{
  enhanced, breakable,
  colback=maflight, colframe=mafamber,
  boxrule=0pt, leftrule=3pt,
  arc=0pt, outer arc=0pt,
  left=8pt, right=8pt, top=6pt, bottom=6pt,
  fonttitle=\bfseries\small,
  title={Version sensitive}, #1
}

% Marks a listing whose exact signature must be re-checked against the SDK.
\newtcolorbox{verifybox}[1][]{
  enhanced, breakable,
  colback=white, colframe=mafgrey,
  boxrule=0.6pt, arc=0pt, outer arc=0pt,
  left=8pt, right=8pt, top=6pt, bottom=6pt,
  fonttitle=\bfseries\small,
  title={Compile before you trust this}, #1
}

\newtcolorbox{exercisebox}[1][]{
  enhanced, breakable,
  colback=white, colframe=mafteal,
  boxrule=0pt, leftrule=3pt,
  arc=0pt, outer arc=0pt,
  left=8pt, right=8pt, top=6pt, bottom=6pt,
  fonttitle=\bfseries\small,
  title={Exercise}, #1
}

% ============================================================
%  SCREENSHOT PLACEHOLDERS
%  ---------------------------------------------------------
%  \needscreenshot{key}{caption}{what exactly to capture}
%
%  If figures/screenshots/<key>.png exists it is included.
%  If not, a visible placeholder box is drawn instead, and the
%  request is recorded in the Screenshot Manifest (Appendix D)
%  via \listofshots.
% ============================================================

\makeatletter
\newcommand{\listofshots}{%
  \section*{Screenshot manifest}%
  \addcontentsline{toc}{section}{Screenshot manifest}%
  \@starttoc{shots}%
}
\makeatother

\newcommand{\needscreenshot}[3]{%
  \addcontentsline{shots}{subsection}{\protect\numberline{}\texttt{#1.png} --- #3}%
  \begin{figure}[htbp]
  \centering
  \IfFileExists{figures/screenshots/#1.png}{%
    \includegraphics[width=0.95\linewidth]{figures/screenshots/#1.png}%
  }{%
    \begin{tcolorbox}[
      enhanced, width=0.95\linewidth,
      colback=maflight, colframe=mafamber,
      boxrule=1pt, arc=0pt, outer arc=0pt,
      borderline={0.6pt}{3pt}{mafamber, dashed},
      left=12pt, right=12pt, top=14pt, bottom=14pt]
      \centering
      {\bfseries\color{mafamber}\large SCREENSHOT NEEDED}\\[6pt]
      {\ttfamily\small figures/screenshots/#1.png}\\[10pt]
      \begin{minipage}{0.9\linewidth}
        \small\raggedright\color{mafgrey} #3
      \end{minipage}
    \end{tcolorbox}%
  }%
  \caption{#2}\label{fig:#1}
  \end{figure}%
}

% A normal (non-screenshot) figure that still needs producing — diagrams etc.
\newcommand{\needdiagram}[3]{%
  \begin{figure}[htbp]
  \centering
  \IfFileExists{figures/#1.pdf}{%
    \includegraphics[width=0.95\linewidth]{figures/#1.pdf}%
  }{%
    \begin{tcolorbox}[enhanced, width=0.95\linewidth,
      colback=white, colframe=mafgrey, boxrule=0.8pt,
      arc=0pt, outer arc=0pt, left=12pt, right=12pt, top=14pt, bottom=14pt]
      \centering{\bfseries\color{mafgrey}DIAGRAM: \ttfamily figures/#1.pdf}\\[8pt]
      \begin{minipage}{0.9\linewidth}\small\raggedright #3\end{minipage}
    \end{tcolorbox}%
  }%
  \caption{#2}\label{fig:#1}
  \end{figure}%
}

% ============================================================
%  PINNED VERSIONS — single source of truth
%  Change these here; they propagate through the whole book.
% ============================================================
\newcommand{\mafcore}{1.10.0}         % core train: Agents.AI, .OpenAI, .Workflows
\newcommand{\mafext}{1.6.2}           % extensions train: Harness, DurableTask, Foundry.Hosting
\newcommand{\mafworkflows}{1.10.0}    % same train as core; kept as a separate macro
\newcommand{\mafdate}{July 2026}      % "verified as of"
\newcommand{\dotnetver}{.NET 10}
\newcommand{\vsver}{Visual Studio 2026}

% The index is two-column and its entries include \texttt{} API names, which do
% not hyphenate. Ragged right keeps multi-word entries off the margin. It cannot
% help a single token wider than the column, so keep verbatim index entries under
% about 29 characters and index longer type names by concept instead.
% Appended to \theindex so it lands after the class's own \twocolumn setup.
\apptocmd{\theindex}{\raggedright}{}{%
  \PackageWarning{maf}{Could not patch theindex for ragged-right setting}}

\makeindex
